% ============================================================================
% 第 2 章 - 用户态 AArch64 汇编
% ============================================================================
\chapter{用户态 AArch64 汇编}

本章我们从零开始，用 AArch64（ARM 64 位）汇编语言在 macOS 上写一个完整的用户态程序。
它没有 \texttt{main} 函数，没有 \texttt{printf}，也没有 \texttt{malloc}——所有这些
我们都用最底层的方式自己实现。

程序的整体结构是：一个 4 KB 的堆内存分配器，一个按序号存储字符串的平行数组，
以及一个命令解析器的雏形。主程序循环工作流程为：
\begin{center}
    读取输入 \(\rightarrow\) 分配内存并复制 \(\rightarrow\)
    存入字符串数组 \(\rightarrow\) 打印 \(\rightarrow\) 扫描空格位置
\end{center}

\section{系统调用：最底层的 I/O}
\input{isa/sections/01_syscall}

\section{堆内存分配器}
\input{isa/sections/02_memory}

\section{字符串复制：用户态的 strcpy}
\input{isa/sections/03_strcopy}

\section{平行数组：字符串的索引存储}
\input{isa/sections/04_string_array}

\section{命令解析：find\_two\_spaces}
\input{isa/sections/05_parse}

\section{主程序：把一切串起来}
% ============================================================================
% 2.6 主程序：把一切串起来
% ============================================================================
\subsection{hello.s：主程序入口}

主程序的工作非常简单——初始化堆之后，就进入一个循环：
读一行→复制→存到平行数组→打印→找空格。

\begin{lstlisting}[style=asm,caption=hello.s — 主程序]
.text
.global _main
.extern exit
.extern memory_init
.extern memory_alloc
.extern memory_free
.extern sys_write
.extern sys_read
.extern test_write
.extern test_memory
.extern test_read
.extern test_write_memory
.extern store_string_ptr_len
.extern get_string_ptr
.extern get_string_len
.extern store_ptr_done
.extern print_index_strings
.extern find_two_spaces

_main:
    bl memory_init             ; 初始化堆
    mov x19, #0               ; x19 = 循环计数器（也是字符串序号）

loop:
    mov x20, #0
    bl store_string           ; 读一行，分配内存，存入数组

    mov x20, #0
    bl get_string_ptr_len_spaces ; 顺便检查空格位置

    add x19, x19, #1
    cmp x19, #1
    b.eq do_exit
    b loop

do_exit:
    b exit
\end{lstlisting}

\subsection{store\_string：读入并保存一行}

\begin{lstlisting}[style=asm,caption=hello.s — store\_string]
store_string:
    stp x29, x30, [sp, #-16]!
    bl test_read              ; 把一行
    bl copy_string            ; copy_string(x0) 从 input_buf 复制到新分配的堆内存
    mov x1, x20                ; 序号
    bl store_string_ptr_len     ; 存入平行数组
    ldp x29, x30, [sp], #16
    ret
\end{lstlisting}

\subsection{get\_string\_ptr\_len\_spaces：打印并解析空格}

\begin{lstlisting}[style=asm,caption=hello.s — 打印+解析]
get_string_ptr_len_spaces:
    stp x29, x30, [sp, #-16]!

    bl print_string           ; 把刚存入的字符串再取出来打印
    stp x1, x2, [sp, #-16]! ; 保存打印返回的指针和长度
    bl find_two_spaces        ; 然后扫描字符串找空格位置
    ldp x2, x3, [sp], #16

get_string_ptr_len_spaces_done:
    ldp x29, x30, [sp], #16
    ret
\end{lstlisting}

\subsection{理解整体数据流}

整个程序的数据流如下：

\begin{verbatim}
stdin → sys_read → input_buf → copy_string(从 input_buf 到 heap)
        → (ptr, len) → store_string_ptr_len
        → print_index_strings(x19) → sys_write → stdout
        → find_two_spaces(buffer) → x0=第一个空格, x1=第二个空格
\end{verbatim}

程序循环一次（\texttt{x19} 递增一，直到 \texttt{x19 = 1} 时退出。
你可以把 \texttt{cmp x19, \#1} 改成更大的数字来延长程序执行更多轮。

\subsection{主循环限制}

这里的循环在 \texttt{x19} 计数，
\textbf{最多能存 MAX\_STRINGS（32）条字符串}。超过会被
\texttt{store\_string\_ptr\_len} 里的 \texttt{cmp x1, \#MAX\_STRINGS; b.hs store\_ptr\_done} 拦截掉——如果 \texttt{x19} 超过 32，函数什么都不做就返回。

\subsection{新版主程序：把解析后的 token 逐条打印}

在引入 \texttt{store\_instr\_addrs} 之后，主程序多了一步：
把刚才那条命令的 4 个 token 分别打印出来，验证解析是否正确。

\begin{lstlisting}[style=asm, caption=hello.s — 新版 \_main（核心循环）]
_main:
    bl memory_init                 ; 初始化堆
    mov x19, #0                    ; x19 = 循环计数器

loop:
    mov x20, #0
    bl  store_string               ; 读一行 → 复制 → 存到平行数组

    mov x20, #0
    bl  get_string_ptr_len_spaces  ; 顺便找两个空格位置，顺便把字符串打印出来
                                   ; 返回后：x0 = 第一个空格, x1 = 第二个空格
                                   ; x2 = 字符串起始, x3 = 字符串长度

    bl  store_instr_addrs          ; 把刚才的一行切成 3 个 token，存进 instr_addrs

    ; ---- 把 token 0/1/2/3 分别打印出来，每个后面换行 ----
    mov x0, #0
    bl  get_instr_addrs            ; x0 = 第 0 条指令的 4 指针数组头
    bl  get_instr_0                ; 打印 token 0
    bl  print_newline

    mov x0, #0
    bl  get_instr_addrs
    bl  get_instr_1                ; 打印 token 1
    bl  print_newline

    mov x0, #0
    bl  get_instr_addrs
    bl  get_instr_2                ; 打印 token 2
    bl  print_newline

    mov x0, #0
    bl  get_instr_addrs
    bl  get_instr_3                ; 打印 token 3
    bl  print_newline

flag:
    add x19, x19, #1
    cmp x19, #1
    b.eq do_exit
    b loop

do_exit:
    b exit
\end{lstlisting}

调用完空格查找函数后，\texttt{x0/x1} 里是两个空格位置，
\texttt{x2/x3} 里是字符串的起始地址和长度，
正好匹配 \texttt{store\_instr\_addrs} 期望的 4 个入参。

\subsection{get\_instr\_0/1/2/3：从指针数组里拿一个 token 打印}

这四个函数的结构几乎一样，只是读取的偏移不同：

\begin{lstlisting}[style=asm, caption=hello.s — get\_instr\_0（打印 token 0）]
get_instr_0:
    stp x29, x30, [sp, #-32]!
    stp x19, x20, [sp, #16]
    mov x20, x0                    ; 保存 4 指针数组头

    ; 手动 strlen: 找到 token 0 的长度
    ldr x1, [x20]                  ; x1 = token 0 的字符串起始
    mov x2, #0
strlen_loop:
    ldrb w3, [x1, x2]
    cbz  w3, strlen_done
    add  x2, x2, #1
    b    strlen_loop
strlen_done:
    mov  x0, #1
    bl   sys_write                 ; sys_write(1, token0_ptr, token0_len)

    ldp x19, x20, [sp, #16]
    ldp x29, x30, [sp], #32
    ret
\end{lstlisting}

\texttt{get\_instr\_1/2/3} 的区别只是把 \texttt{ldr x1, [x20]}
换成 \texttt{ldr x1, [x20, \#8]}、\texttt{ldr x1, [x20, \#16]}、
\texttt{ldr x1, [x20, \#24]}——分别取 token 1、2、3 的指针。

\subsection{print\_newline：输出一个换行符}

在两个 token 之间打印 \texttt{\textbackslash n}，让每个 token 单独占一行：

\begin{lstlisting}[style=asm, caption=hello.s — print\_newline]
print_newline:
    stp  x29, x30, [sp, #-32]!
    strb wzr, [sp, #16]            ; 清一块栈内存
    mov  w8, #10                   ; '\n' 的 ASCII = 10
    strb w8,  [sp, #16]            ; 把换行符写进去
    add  x1,  sp, #16              ; x1 = 换行符所在地址
    mov  x0,  #1
    mov  x2,  #1
    bl   sys_write                 ; sys_write(1, "\n", 1)
    ldp  x29, x30, [sp], #32
    ret
\end{lstlisting}

它用栈上的一小块空间临时存放换行符，不需要专门准备一个 \texttt{.data}
里的字符串。这是"不依赖静态数据、临时在栈上构造参数"的常见小技巧。

\subsection{现在整个程序的工作流}

综合起来，一次循环做了以下 4 件事：

\begin{enumerate}
    \item 从 \code{stdin} 读一行 → \code{copy\_string} 复制到堆 →
        把指针和长度存入 \code{string\_ptrs} 与 \code{string\_lens} 两个数组
    \item 打印那一行 → 扫描字符串找两个空格位置
    \item 用空格位置把一行切成 3 个 token → 存入 \code{instr\_addrs}
    \item 按序号把 4 个 token 分别打印出来
\end{enumerate}

这相当于一个最小的"命令解析 + 回显"循环：
你输入什么，它就原样打印，再把每个 token 拆出来单独打印一行。


\section{编译与调试脚本}
\input{isa/sections/07_build}

\section{指令解析：把命令切成 token 指针数组}
% ============================================================================
% 2.8 指令解析：把空格位置变成 token 指针数组
% ============================================================================

\texttt{find\_two\_spaces} 告诉我们"第一个空格在哪、第二个空格在哪"，
但要真正执行命令，我们还需要把\textbf{各个 token 的指针保存下来}，
后面按序号取回、并做指令比对。这就是 \texttt{instr.s} 要做的事：

\begin{enumerate}
    \item 分配一块内存，里面放 4 个 8 字节指针（最多 4 个 token）
    \item 把 [0, 第一个空格)、[第一个空格+1, 第二个空格)、[第二个空格+1, 末尾)
        三段分别复制到新分配的内存中（每段末尾补 0）
    \item 把这 4 个指针的\textbf{地址}放到全局的 \texttt{instr\_addrs} 数组里，
        按序号保存，供后续读取
\end{enumerate}

\subsection{数据区：instr\_addrs 全局指针数组}

\texttt{instr.s} 开头定义了一个与 \texttt{string\_array.s} 类似的平行数组，
区别是每项是一个"指针数组的头"（也就是一块包含 4 个指针的内存）：

\begin{lstlisting}[style=asm, caption=instr.s — 数据区定义]
.include "src/memory.inc"

.data
.align 4
.extern heap
.extern heap_end
.equ MAX_INSTR, 64             ; 最多保存 64 条已解析的指令
instr_addrs: .space MAX_INSTR*8  ; 每项 8 字节：指向一条 token 数组的指针
count_instr: .word 0             ; 当前已存入多少条

.text
.extern memory_init
.extern memory_alloc
.extern memory_free

.global store_instr_addrs
.global get_instr_addrs
\end{lstlisting}

布局上，\texttt{instr\_addrs} 是一个指针数组：
第 0 项指向第 0 条指令的 4 个 token 指针块，
第 1 项指向第 1 条指令的 4 个 token 指针块，以此类推。

每条指令的"4 个 token 指针块"是这样一块内存：

\begin{center}
\begin{tabular}{|l|l|}
\hline
偏移 0  & token\_0 的指针（如 "add" 的首地址） \\
\hline
偏移 8  & token\_1 的指针（如 "file1" 的首地址） \\
\hline
偏移 16 & token\_2 的指针（如 "file2" 的首地址） \\
\hline
偏移 24 & token\_3 的指针（保留 / 未使用） \\
\hline
\end{tabular}
\end{center}

\subsection{store\_instr\_addrs：把一条命令切成 token 并存下}

它接收 4 个参数：
\begin{itemize}
    \item \texttt{x0} = 第一个空格位置
    \item \texttt{x1} = 第二个空格位置
    \item \texttt{x2} = 原字符串的起始地址
    \item \texttt{x3} = 原字符串的长度
\end{itemize}

首先检查计数器，然后分配一块 32 字节（4 × 8）的内存用于存放 4 个 token 指针，
接着分 3 次把每个 token 复制到新分配的独立内存里，
每段末尾补 0 做 C 风格字符串结尾。最后把这块 4 指针的地址存入全局的
\texttt{instr\_addrs[count\_instr]}。

\begin{lstlisting}[style=asm, caption=instr.s — store\_instr\_addrs（入栈与检查）]
store_instr_addrs:
    stp x19, x20, [sp, #-96]!
    stp x21, x22, [sp, #16]
    stp x23, x24, [sp, #32]
    stp x25, x26, [sp, #48]
    stp x27, x28, [sp, #64]
    stp x29, x30, [sp, #80]

    mov x19, x0                 ; x19 = 第一个空格位置
    mov x20, x1                 ; x20 = 第二个空格位置
    mov x21, x2                 ; x21 = 字符串起始地址
    mov x22, x3                 ; x22 = 字符串总长度

    adrp x25, count_instr@PAGE
    add  x25, x25, count_instr@PAGEOFF
    ldr  w26, [x25]             ; 读当前计数
    cmp  w26, #MAX_INSTR
    b.hs skip_store             ; 已达上限，跳过

    mov  x0, #32
    bl   memory_alloc           ; 分配 32 字节 = 4 个指针
    mov  x24, x0                ; x24 = token 指针数组的头

    adrp x23, instr_addrs@PAGE
    add  x23, x23, instr_addrs@PAGEOFF
\end{lstlisting}

接下来三段几乎对称的代码分别把 token\_0、token\_1、token\_2 复制出来：

\begin{lstlisting}[style=asm, caption=instr.s — 解析并复制 3 个 token]
; ---- token 0: [0, x19) ----
parse_instr1:
    mov x0, x19                 ; 长度 = 第一个空格位置
    add x0, x0, #1              ; +1 给结尾 0 留位置
    bl  memory_alloc            ; 分配新内存存放这一段
    mov x27, x0                 ; x27 = 目标地址
    mov x28, #0
extract_byte_loop1:
    cmp x28, x19
    b.ge extract_done1
    ldrb w29, [x21, x28]
    strb w29, [x27, x28]
    add x28, x28, #1
    b extract_byte_loop1
extract_done1:
    strb wzr, [x27, x28]        ; 末尾补 0
    str  x27, [x24]             ; token 0 指针写入 instr_addrs[0]

; ---- token 1: (x19, x20) ----
parse_instr2:
    sub x0, x20, x19            ; 长度 = 第二个空格 - 第一个空格
    add x0, x0, #1
    bl  memory_alloc
    mov x27, x0
    mov x28, #0
    add x29, x19, #1            ; 从第一个空格之后开始读
extract_byte_loop2:
    cmp x29, x20
    b.ge extract_done2
    ldrb w30, [x21, x29]
    strb w30, [x27, x28]
    add x28, x28, #1
    add x29, x29, #1
    b extract_byte_loop2
extract_done2:
    strb wzr, [x27, x28]
    str  x27, [x24, #8]         ; token 1 指针写入 instr_addrs[1]

; ---- token 2: (x20, end) ----
parse_instr3:
    sub x0, x22, x20            ; 长度 = 总长 - 第二个空格
    add x0, x0, #1
    bl  memory_alloc
    mov x27, x0
    mov x28, #0
    add x29, x20, #1            ; 从第二个空格之后开始
extract_byte_loop3:
    cmp x29, x22
    b.ge extract_done3
    ldrb w30, [x21, x29]
    strb w30, [x27, x28]
    add x28, x28, #1
    add x29, x29, #1
    b extract_byte_loop3
extract_done3:
    strb wzr, [x27, x28]
    str  x27, [x24, #16]        ; token 2 指针写入 instr_addrs[2]

    str  x27, [x24, #24]        ; token 3 保留（与 token 2 同值，方便调试）
\end{lstlisting}

最后把这块 4 指针数组的头地址写入全局 \texttt{instr\_addrs}，计数器加一：

\begin{lstlisting}[style=asm, caption=instr.s — 收尾：写入全局指针数组]
    str  x24, [x23, x26, lsl #3]  ; instr_addrs[count_instr] = x24

    mov  w0, w26                   ; 返回值 = 这条指令的序号
    add  w26, w26, #1
    str  w26, [x25]                ; count_instr++

instr_addrs_done:
    ldp x29, x30, [sp, #80]
    ldp x27, x28, [sp, #64]
    ldp x25, x26, [sp, #48]
    ldp x23, x24, [sp, #32]
    ldp x21, x22, [sp, #16]
    ldp x19, x20, [sp], #96
    ret

skip_store:
    mov x0, #-1                    ; 失败返回 -1
    b instr_addrs_done
\end{lstlisting}

\subsection{get\_instr\_addrs：按序号取回 token 指针数组}

保存之后，我们希望按"这是第几条指令"把 4 个 token 的指针数组读回来。
这是写入的反向操作：

\begin{lstlisting}[style=asm, caption=instr.s — get\_instr\_addrs]
get_instr_addrs:
    stp x29, x30, [sp, #-32]!
    stp x27, x28, [sp, #16]

    mov  x28, x0                   ; x28 = 要读取的指令序号
    adrp x29, instr_addrs@PAGE
    add  x29, x29, instr_addrs@PAGEOFF
    ldr  x0, [x29, x28, lsl #3]    ; x0 = instr_addrs[序号]

get_instr_addrs_done:
    ldp x27, x28, [sp, #16]
    ldp x29, x30, [sp], #32
    ret
\end{lstlisting}

调用它之后，\texttt{x0} 就是一块 32 字节的内存，
里面依次放着 token\_0、token\_1、token\_2、token\_3 的字符串指针。
主程序里的 \texttt{get\_instr\_0 / get\_instr\_1 / get\_instr\_2 / get\_instr\_3}
就是把 \texttt{x0} 当作"指针数组"来用的例子。

\subsection{为什么要单独复制 token？}

一个很自然的问题：\emph{既然原字符串已经在堆里了，
为什么不直接保存"偏移 + 长度"，而是额外分配内存把每个 token 复制一份？}

原因有两个：

\begin{enumerate}
    \item \textbf{C 风格字符串方便后续比对}。后面如果想写
        \texttt{strcmp("add", token0)} 这种比较，要求被比较的字符串
        在自己的末尾处有一个 \texttt{0} 字节。原字符串在 token 交界处
        是空格而不是 \texttt{0}，所以必须复制一份并补 \texttt{0}。
    \item \textbf{后续可能修改或释放原字符串}。一旦原字符串所在的行
        被释放或被覆盖，我们保存在 \texttt{instr\_addrs} 里的"指向原字符串
        内部的指针"就会变成野指针。复制之后 token 各自独立，不再依赖原件。
\end{enumerate}

代价是每条指令会多 3 次 \texttt{memory\_alloc}，
但我们的堆有 4096 字节，一个 token 通常只有几个字符，在实验阶段完全够用。

